\documentclass[11pt,a4paper]{article}
\usepackage[utf8]{inputenc}\usepackage[T1]{fontenc}\usepackage{lmodern}
\usepackage{amsmath,amssymb,amsthm}\usepackage{geometry}\usepackage{xcolor}
\usepackage{hyperref}\usepackage{microtype}\usepackage{enumitem}
\geometry{margin=2.5cm}
\hypersetup{colorlinks=true,linkcolor=blue!55!black,urlcolor=blue!55!black}
\theoremstyle{plain}
\newtheorem{theorem}{Theorem}[section]
\newtheorem{lemma}[theorem]{Lemma}
\newtheorem{corollary}[theorem]{Corollary}
\theoremstyle{definition}
\newtheorem{remark}[theorem]{Remark}
\newcommand{\CC}{C_{\mathrm{Collatz}}}

\title{\bfseries Transcendence of a Collatz-Related Constant\\ via Sturmian Expansions}
\author{Joel Cruz Cabrera\\
\small Santa Cruz de Tenerife, Canary Islands, Spain\\
\small\texttt{tf.joel.work@gmail.com}\quad ORCID: 0009-0005-4048-1237}
\date{July 2026}

\sloppy
\begin{document}\maketitle

\begin{abstract}
Let $\CC=\tfrac12-\sum_{n\ge2}2^{-\lfloor n\log_2 3\rfloor}=0.2862745023241087\ldots$, a
constant arising as a natural density in the analysis of the accelerated Collatz map. We show
that the binary expansion of $\CC$ is a Sturmian word, and deduce that $\CC$ is
transcendental. The argument is not specific to the base $3$: for every $a$ with $\log_2 a$
irrational, the constant $C_a=1-\sum_{n\ge1}2^{-\lfloor n\log_2 a\rfloor}$ is transcendental.
\textbf{Nothing here bears on the truth of the Collatz conjecture}; the result is a statement
in transcendence theory that happens to concern a constant met in that context.
\end{abstract}

\noindent\textbf{MSC:} 11J81 (Transcendence), 68R15 (Combinatorics on words), 11B83.

\section{The constant}
For odd $n$, the accelerated Collatz step is $\operatorname{step}(x)=(3x+1)/2^{v_2(3x+1)}$.
Counting the odd integers whose first block does not immediately shrink leads to the density
\begin{equation}\label{eq:def}
  \CC \;=\; \tfrac12-\sum_{n\ge2}2^{-\lfloor n\log_2 3\rfloor}.
\end{equation}
Summing \eqref{eq:def} to $n=300$ in $80$-digit arithmetic gives
\[
  \CC=0.28627450232410874196\ldots
\]
(the tail beyond $n=300$ is below $2^{-475}$). An equivalent closed form,
$1-\sum_{n\ge1}2^{-(n+\lfloor n\log_2(3/2)\rfloor)}=\CC$, is machine-checked in Lean~4; it
follows from the elementary identity $n+\lfloor n\log_2(3/2)\rfloor=\lfloor n\log_2 3\rfloor$
(a one-line consequence of $\lfloor x-k\rfloor=\lfloor x\rfloor-k$ for $k\in\mathbb{Z}$).

\section{The expansion is Sturmian}
Write $\beta=\log_2 3\in(1,2)$, which is irrational, and let
$B=(\lfloor n\beta\rfloor)_{n\ge1}$ be the associated Beatty sequence. Define the gap word
$w=(w_n)_{n\ge1}$ by $w_n=\lfloor n\beta\rfloor-\lfloor (n-1)\beta\rfloor$.

\begin{lemma}[Two-letter alphabet]\label{lem:gap}
$w_n\in\{1,2\}$ for all $n$.
\end{lemma}
\begin{proof}
$\lfloor x+\beta\rfloor-\lfloor x\rfloor\in\{\lfloor\beta\rfloor,\lfloor\beta\rfloor+1\}
=\{1,2\}$ for any $x$, since $1<\beta<2$.
\end{proof}

\begin{lemma}[Sturmian]\label{lem:sturm}
$w$ is a Sturmian word: its factor complexity satisfies $p(k)=k+1$ for every $k\ge1$.
\end{lemma}
\begin{proof}
$w$ is, by Lemma~\ref{lem:gap}, the difference word of a Beatty sequence of irrational slope
$\beta$; such words are exactly the mechanical words of slope $\beta-\lfloor\beta\rfloor$,
which are Sturmian precisely because the slope is irrational (Morse--Hedlund~\cite{mh}).
\end{proof}
\noindent \emph{Verification.} Computing the factors of $w$ directly for $n\le4000$ gives
$p(k)=k+1$ for all $k\le8$, with alphabet $\{1,2\}$, as predicted.

Because the exponents $\lfloor n\beta\rfloor$ increase by $w_n\in\{1,2\}$, the binary digits
of the series in \eqref{eq:def} are supported exactly on the Beatty set $\{\lfloor
n\beta\rfloor\}$; the digit word of $\CC$ is therefore the characteristic word of that set,
which is Sturmian by Lemma~\ref{lem:sturm}.

\section{Transcendence}
\begin{theorem}\label{thm:main}
$\CC$ is transcendental.
\end{theorem}
\begin{proof}
By Section~2 the binary expansion of $\CC$ is a Sturmian word, hence has factor complexity
$p(k)=k+1$, which is $o(k\log k)$ and in particular sub-linear in the sense required.
A real number whose base-$b$ expansion is Sturmian is transcendental: this is the theorem of
Ferenczi--Mauduit~\cite{fm}, and follows also from the complexity criterion of
Adamczewski--Bugeaud~\cite{ab}, since an algebraic irrational cannot have an expansion of
complexity $p(k)=k+1$. As $\CC$ is irrational (its expansion is aperiodic, being Sturmian),
it is transcendental.
\end{proof}

\begin{corollary}
For every real $a>1$ with $\log_2 a\notin\mathbb{Q}$, the constant
$C_a=1-\sum_{n\ge1}2^{-\lfloor n\log_2 a\rfloor}$ is transcendental. In particular this holds
for every integer $a$ that is not a power of~$2$.
\end{corollary}
\begin{proof}
Identical, with $\beta=\log_2 a$; only irrationality of $\beta$ is used. If $1<\beta<2$
fails, replace $w$ by the gap word of $\lfloor n\beta\rfloor$, whose alphabet is
$\{\lfloor\beta\rfloor,\lfloor\beta\rfloor+1\}$; the mechanical-word argument is unchanged.
\end{proof}

\begin{remark}
We checked $\CC$ against the CODATA physical constants and the standard mathematical
constants and found no match to within $10^{-3}$; we know of no closed form for it, and it is transcendental by Theorem~\ref{thm:main}. It is
not, however, a constant of a new kind: it belongs to the classical family
$\sum_{n\ge1} b^{-\lfloor n\alpha\rfloor}$, studied since B\"ohmer~\cite{bohmer}, who proved
transcendence when $\alpha$ has unbounded partial quotients, and later by
Danilov~\cite{danilov}, Davison~\cite{davison} and Adams--Davison~\cite{adamsdavison}, the last
of whom give the continued fraction of $(b-1)\sum_{n\ge1} b^{-\lfloor n\alpha\rfloor}$
explicitly for every irrational $\alpha$. What is specific here is only that this member of
the family, with $b=2$ and $\alpha=\log_2 3$, is the one the Collatz observable produces.
\end{remark}

\begin{remark}[Diophantine flavour]
The continued fraction of $\CC$ begins
\[
  [0;\,3,\,2,\,36,\,2056,\,70368878395648,\,\dots];
\]
the very large early partial quotients reflect exceptionally good dyadic-rational
approximations. Only the terms stable across two independent precisions are reported.
\end{remark}

\section*{Reproducibility}
The value, the gap alphabet and the factor complexity are recomputed by
\texttt{tools/verify\_for\_papers.py}; the closed-form identity is checked by Lean~4
(\texttt{blockSum\_eq\_cCollatzClosed}, no custom axioms).

\section*{Acknowledgement}
AI-assisted; every claim was verified by exact computation or Lean's kernel before inclusion.

\begin{thebibliography}{99}
\bibitem{fm} S.~Ferenczi, C.~Mauduit, \emph{Transcendence of numbers with a low complexity
expansion}, J. Number Theory \textbf{67} (1997), 146--161.
\bibitem{ab} B.~Adamczewski, Y.~Bugeaud, \emph{On the complexity of algebraic numbers I.
Expansions in integer bases}, Ann. of Math. \textbf{165} (2007), 547--565.
\bibitem{lothaire} M.~Lothaire, \emph{Algebraic Combinatorics on Words}, CUP, 2002.
\bibitem{lagarias} J.~C. Lagarias, \emph{The $3x+1$ problem and its generalizations},
Amer. Math. Monthly \textbf{92} (1985), 3--23.
\bibitem{mh} M.~Morse and G.~A. Hedlund, \emph{Symbolic dynamics II: Sturmian trajectories},
Amer. J. Math. \textbf{62} (1940), 1--42. doi:10.2307/2371431.
\bibitem{bohmer} P.~E. B\"ohmer, \emph{\"Uber die Transzendenz gewisser dyadischer Br\"uche},
Math. Ann. \textbf{96} (1927), 367--377. doi:10.1007/BF01209172.
\bibitem{danilov} L.~V. Danilov, \emph{Some classes of transcendental numbers},
Math. Notes \textbf{12} (1972), 524--527. doi:10.1007/BF01095010.
\bibitem{davison} J.~L. Davison, \emph{A series and its associated continued fraction},
Proc. Amer. Math. Soc. \textbf{63} (1977), 29--32. doi:10.1090/S0002-9939-1977-0429778-5.
\bibitem{adamsdavison} W.~W. Adams and J.~L. Davison, \emph{A remarkable class of continued
fractions}, Proc. Amer. Math. Soc. \textbf{65} (1977), 194--198.
doi:10.1090/S0002-9939-1977-0441879-4.
\end{thebibliography}
\end{document}
